\documentclass[a4paper]{article}

% Paquetes básicos
% Español (México) con XeLaTeX: fontspec para fuentes Unicode, polyglossia
% para la división silábica y los nombres del documento en la variante mexicana.
\usepackage{fontspec}
\usepackage{polyglossia}
\setdefaultlanguage[variant=mexican]{spanish}
% Los nombres chinos van como «pinyin (original)»: xeCJK da a los caracteres
% Han su propia fuente; sin ella XeLaTeX los omite con un aviso.
% FandolSong no cubre algunos caracteres de nombres (U+73FA); WenQuanYi Zen Hei sí.
\usepackage[AutoFallBack=true]{xeCJK}
\setCJKmainfont{FandolSong-Regular.otf}
\setCJKfallbackfamilyfont{\CJKrmdefault}{WenQuanYi Zen Hei}
% polyglossia no traduce los nombres de listados.
\AtBeginDocument{\providecommand{\lstlistingname}{}\renewcommand{\lstlistingname}{Listado}%
  \providecommand{\lstlistlistingname}{}\renewcommand{\lstlistlistingname}{Índice de listados}}
\usepackage{amsmath, amssymb}
\usepackage{graphicx}
\usepackage[margin=2.5cm]{geometry}
\usepackage[most]{tcolorbox}
\usepackage{etoolbox}
\usepackage{listings}
\usepackage{hyperref}
\usepackage{booktabs}
\usepackage{subcaption}
\usepackage{float}

% Cajas de resaltado
\newtcolorbox{knowledgebox}[1]{
    enhanced, colback=blue!5!white, colframe=blue!75!black, colbacktitle=blue!75!black,
    coltitle=white, fonttitle=\bfseries, title={#1}, attach boxed title to top left={yshift=-2mm, xshift=2mm},
    boxrule=1pt, sharp corners
}

\newtcolorbox{importantbox}[1]{
    enhanced, colback=yellow!10!white, colframe=yellow!80!black, colbacktitle=yellow!80!black,
    coltitle=black, fonttitle=\bfseries, title={#1}, sharp corners
}

\newtcolorbox{warningbox}[1]{
    enhanced, colback=red!5!white, colframe=red!75!black, colbacktitle=red!75!black,
    coltitle=white, fonttitle=\bfseries, title={#1}, sharp corners
}

% Listados
\lstset{
    language=Python,
    basicstyle=\ttfamily\small,
    keywordstyle=\color{blue},
    stringstyle=\color{red!60!black},
    commentstyle=\color{green!60!black},
    breaklines=true,
    frame=single,
    numbers=left,
    numberstyle=\tiny\color{gray},
    captionpos=b,
    extendedchars=true
}

% Metadatos de la portada
\newcommand{\notetitle}{Clase 11 del curso KAIST CS492D: Generación guiada en tiempo de inferencia 1}
\newcommand{\noteauthors}{Elaboradas a partir de la clase de Minhyuk Sung}
\newcommand{\notedate}{Otoño de 2025}
\newcommand{\videochannel}{Minhyuk Sung (KAIST)}
\newcommand{\videopublishdate}{2025}
\newcommand{\videoduration}{58 minutos}
\newcommand{\videourl}{https://www.youtube.com/watch?v=AjrLJeWKf3M}
\newcommand{\videocoverpath}{cover.jpg}
\newcommand{\repourl}{https://github.com/hqhq1025/ai-course-notes}

% Pie de página con información del repositorio
\usepackage{fancyhdr}
\pagestyle{fancy}
\fancyhf{}
\fancyfoot[C]{\thepage}
\fancyfoot[R]{\footnotesize\color{gray}\href{\repourl}{AI Course Notes}}
\renewcommand{\headrulewidth}{0pt}
\renewcommand{\footrulewidth}{0pt}

\begin{document}

````

\begin{titlepage}
\centering
{\Large Notas del curso\par}
\vspace{1.2cm}
{\huge\bfseries \notetitle\par}
\vspace{0.8cm}
{\large \noteauthors\par}
\vspace{0.3cm}
{\large \notedate\par}
\vspace{1.2cm}

\ifdefempty{\videocoverpath}{
{\small Al generar las notas, indique la ruta local de la portada del video.\par}
}{
\includegraphics[width=0.82\textwidth,height=0.45\textheight,keepaspectratio]{\videocoverpath}\par
}

\vfill
\begin{tcolorbox}[width=0.9\textwidth, colback=black!2!white, colframe=black!60, sharp corners]
\textbf{Autor o canal del video}: \videochannel\par
\textbf{Fecha de publicación}: \videopublishdate\par
\textbf{Duración del video}: \videoduration\par
\ifdefempty{\videourl}{}{
\textbf{Enlace del video}: \href{\videourl}{\nolinkurl{\videourl}}\par
}\par
\textbf{Página del proyecto}: \href{\repourl}{\nolinkurl{\repourl}}\par
\end{tcolorbox}
\end{titlepage}

\tableofcontents
\newpage

%% --- Inicio del contenido --- %%

\section{Introducción: motivación para la guía en tiempo de inferencia}

En las semanas anteriores del curso, discutimos los procesos básicos de entrenamiento y generación de modelos de difusión (Diffusion Model) y modelos de flujo (Flow Model), así como cómo acelerar la generación mediante técnicas en tiempo de inferencia o ajuste fino. El problema central de esta clase gira en torno a otra dirección: \textbf{¿cómo aprovechar los modelos preentrenados para inyectar señales de guía durante la inferencia, de modo que los resultados generados cumplan preferencias o condiciones específicas del usuario?}

\begin{importantbox}{Problema central de la lección}
Dado un modelo de difusión/flujo preentrenado (sin condiciones), ¿cómo podemos guiar el proceso de generación en tiempo de inferencia sin volver a entrenar, para que los puntos de datos resultantes satisfagan condiciones específicas —por ejemplo, coincidan con observaciones parciales dadas, sigan ciertas restricciones de disposición o de estilo?
\end{importantbox}

Actualmente existen numerosos modelos de difusión preentrenados que abarcan múltiples modalidades, como generación de imágenes, generación de video, diseño molecular y planificación del movimiento robótico. Estos modelos suelen requerir grandes volúmenes de datos y capacidad computacional para su entrenamiento, por lo que es difícil para investigadores comunes entrenarlos desde cero. Por ello, surge una pregunta natural: \textbf{¿es posible hacer más cosas utilizando únicamente los modelos preentrenados?}

Una ventaja clave de los modelos de difusión/flujo radica en su característica de \textbf{generación iterativa}. A diferencia de los modelos de generación de un solo paso como GAN y VAE, en los modelos de difusión, a lo largo de la trayectoria desde el ruido $x_T$ hasta los datos limpios $x_0$, cada paso puede ser "manipulado": podemos modificar la velocidad predicha o el ruido para curvar la trayectoria hacia el punto final deseado.

\begin{knowledgebox}{Repaso: Classifier Guidance y Classifier-Free Guidance}
En clases anteriores, ya hemos visto dos métodos para inyectar guía durante el proceso de generación:
\begin{itemize}
    \item \textbf{Classifier Guidance}: Entrena una red neuronal adicional de clasificación que utiliza sus gradientes sobre los datos de ruido intermedios $x_t$ para proporcionar un puntaje de verosimilitud (likelihood score), ajustando así la trayectoria de generación.
    \item \textbf{Classifier-Free Guidance (CFG)}: Utiliza la misma red para aprender simultáneamente el ruido condicional y el incondicional, calculando implícitamente el puntaje de verosimilitud a través de su diferencia y controlando la intensidad de la señal de guía mediante un parámetro $w$.
\end{itemize}
Ambos métodos requieren \textbf{entrenamiento}: ya sea entrenar un clasificador adicional o incluir pares de datos condicionales e incondicionales durante el entrenamiento del predictor de ruido. Esta lección aborda métodos de guía en tiempo de inferencia que \textbf{no requieren entrenamiento adicional}.
\end{knowledgebox}

\subsection{Tres categorías principales de guía en tiempo de inferencia}

El ponente clasifica las técnicas de guía en tiempo de inferencia en tres grandes grupos:

\begin{enumerate}
    \item \textbf{Score Manipulation (manipulación de puntajes)}: Modificaciones prácticas directas a la función de puntajes durante el proceso de desruido, sin garantías teóricas estrictas, pero con buenos resultados en la práctica.
    \item \textbf{Guía basada en marcos de problemas inversos}: Modela la generación condicional como un problema inverso (Inverse Problem), utilizando la fórmula de Bayes para aproximar el cálculo del puntaje de verosimilitud, lo que otorga cierta base teórica.
    \item \textbf{Otras técnicas avanzadas}: Se discutirán en la próxima lección (Lecture 12).
\end{enumerate}

Esta lección cubre principalmente las primeras dos categorías de métodos.

\subsection{Resumen de la sección}

La generación guiada en tiempo de inferencia es una ventaja central de los modelos de difusión, derivada de su naturaleza de generación iterativa. En esta lección, partiremos desde técnicas prácticas de manipulación de puntajes para adentrarnos progresivamente en marcos de resolución de problemas inversos con soporte teórico.



\section{Manipulación de puntuaciones: técnicas prácticas de manipulación de puntajes}

El primer enfoque consiste en manipular directamente la función de puntuación durante el proceso de eliminación de ruido. Este tipo de métodos no cuenta con una explicación teórica estricta, pero suelen funcionar bien en la práctica, habiendo generado numerosos casos de éxito en los campos de imagen y video.

\subsection{RePaint: un caso clásico de reparación de imágenes}

\textbf{RePaint} es un ejemplo típico de manipulación de puntuaciones; resuelve el problema de \textbf{reparación de imágenes (Inpainting)}: dado una imagen parcialmente oculta, rellenar las zonas faltantes para que el resultado sea natural y coherente con el fondo conocido.

\begin{importantbox}{La idea central de RePaint}
En cada paso de eliminación de ruido, se divide el punto de datos intermedio $x_t$ en dos partes para procesarlas:
\begin{itemize}
    \item \textbf{Zona de primer plano} (zona por rellenar): ejecutar el \textbf{proceso inverso} (reverse process) para generar contenido gradualmente a partir del ruido.
    \item \textbf{Zona de fondo} (zona conocida): ejecutar el \textbf{proceso hacia adelante} (forward process), añadiendo ruido a la imagen limpia conocida hasta el nivel de ruido correspondiente al paso temporal $t-1$.
\end{itemize}
Luego se unen ambas partes para obtener $x_{t-1}$, repitiéndose este proceso hasta que $t=0$.
\end{importantbox}

El flujo del algoritmo específico es el siguiente:
\begin{enumerate}
    \item Dado el punto de datos intermedio actual $x_t$, ejecutar un paso de proceso inverso para obtener la zona de primer plano $x_{t-1}^{\text{fg}}$.
    \item Ejecutar un paso de proceso hacia adelante en la imagen de fondo limpia conocida hasta el paso temporal $t-1$, obteniendo $x_{t-1}^{\text{bg}}$.
    \item Combinar ambos utilizando una máscara (mask): tomar $x_{t-1}^{\text{fg}}$ para el primer plano y $x_{t-1}^{\text{bg}}$ para el fondo.
    \item Repetir el proceso anterior.
\end{enumerate}

\begin{warningbox}{RePaint carece de garantías teóricas}
Por qué se conserva el fondo es intuitivo: en cada paso, añadimos ruido de nuevo a partir de la imagen limpia. Sin embargo, no existe garantía teórica sobre la \textbf{integración perfecta} (seamless integration) entre primer plano y fondo. En uso real, efectivamente existen casos fallidos donde el primer plano y el fondo no coinciden. El docente enfatiza: ``No hay teoría aquí, solo ideas prácticas.''
\end{warningbox}

\subsection{Manipulación de mapas de atención}

Además de manipular directamente la puntuación/ruido, también se puede lograr una guía mediante la manipulación de los \textbf{mapas de atención} (Attention Maps) dentro de la red predictora de ruido, incluyendo:

\begin{itemize}
    \item \textbf{Manipulación de auto-atención}: modificar los pesos de atención en el módulo de atención automática para influir en las asociaciones entre diferentes regiones de la imagen.
    \item \textbf{Manipulación de atención cruzada}: modificar la correspondencia entre los tokens de texto y las regiones de imagen en la atención cruzada, controlando así la posición y las atributos de objetos específicos.
\end{itemize}

\begin{knowledgebox}{Ejemplos de aplicación de la manipulación de atención cruzada}
Mediante la manipulación de la atención cruzada, se pueden lograr las siguientes tareas (todas sin necesidad de volver a entrenar el modelo de difusión):
\begin{itemize}
    \item \textbf{Transferencia de estilo}: extraer información estructural de una imagen y información de estilo/apariencia de otra para generar una nueva imagen.
    \item \textbf{Control de disposición}: dado un texto descriptivo y cajas delimitadoras 2D, controlar la posición de cada objeto en la imagen. Por ejemplo, especificar que el ``coche amarillo'' esté dentro del recuadro amarillo a la derecha y el ``castillo'' dentro del recuadro verde a la izquierda.
    \item \textbf{Edición de video}: mediante inversión (inversion) y manipulación de atención, mover objetos en un video a nuevas posiciones mientras se mantiene la naturalidad del video.
\end{itemize}
\end{knowledgebox}

\subsection{Ventajas y desventajas de la manipulación de puntuaciones}

\begin{table}[H]
\centering
\begin{tabular}{p{0.45\textwidth}p{0.45\textwidth}}
\toprule
\textbf{Ventajas} & \textbf{Desventajas} \\
\midrule
Directo, intuitivo y fácil de implementar & Sin garantías teóricas \\
Numerosos casos de éxito en los campos de imagen/video & Existen casos fallidos; los resultados son impredecibles \\
No requiere entrenamiento adicional ni datos & Mayormente métodos ad-hoc diseñados para tareas específicas \\
Extensible a otras modalidades de datos & Difícil de generalizar a nuevos tipos de guía \\
\bottomrule
\end{tabular}
\caption{Comparación de ventajas y desventajas de los métodos de manipulación de puntuaciones}
\end{table}

\subsection{Resumen de la sección}

La manipulación de puntuaciones es una clase de técnicas de guía en tiempo de inferencia altamente prácticas. Su idea central consiste en modificar directamente la función de puntuación o los mapas de atención en cada paso del proceso de eliminación de ruido, haciendo que la trayectoria de generación se incline hacia la dirección deseada. Aunque carece de garantías teóricas, ya existen numerosos casos de éxito en tareas como reparación de imágenes, transferencia de estilo, control de disposición y edición de video. A continuación discutiremos una clase de métodos más sistemáticos y con base teórica: la guía basada en problemas inversos.

\section{Marco de problemas inversos}

La idea central de la segunda categoría de técnicas de guía es modelar la generación condicional como un \textbf{problema inverso}. Se trata de un marco más estructurado y general, aplicable a una gama más amplia de escenarios.

\subsection{Definición del problema inverso}

\begin{importantbox}{Definición formal del problema inverso}
Dado:
\begin{itemize}
    \item Una función de mapeo hacia adelante \textbf{forward mapping function} $\mathcal{A}$, que mapea datos limpios $x$ a observaciones parciales $y$: $y = \mathcal{A}(x) + n$, donde $n$ es el posible ruido de observación.
    \item Un valor de observación $y$.
\end{itemize}
Objetivo: encontrar un $x$ razonable tal que $\mathcal{A}(x) \approx y$.

Dado que $\mathcal{A}$ suele ser \textbf{no invertible}, pueden existir múltiples $x$ razonables para un $y$ dado; por lo tanto, esto es esencialmente un problema de generación condicional.
\end{importantbox}

\subsection{Aplicaciones típicas del problema inverso}

Los problemas inversos tienen numerosas aplicaciones en el procesamiento de imágenes y pueden extenderse a otras modalidades de datos:

\begin{table}[H]
\centering
\begin{tabular}{lll}
\toprule
\textbf{Tarea} & \textbf{Mapeo hacia adelante $\mathcal{A}$} & \textbf{Descripción} \\
\midrule
Denoising gaussiano (desenfoque) & Núcleo de convolución gaussiano & Recuperar una imagen nítida desde una borrosa \\
Superresolución & Operación de submuestreo & Recuperar alta resolución desde baja resolución \\
Inpainting de imágenes & Operación de máscara & Recuperar una imagen completa desde observaciones parciales \\
Reconstrucción CT/MRI & Operación de proyección & Reconstruir 3D desde proyecciones 2D multángulo \\
Colorización de imágenes & Operación de descoloración & Recuperar una imagen a color desde un mapa de grises \\
\bottomrule
\end{tabular}
\caption{Escenarios típicos de aplicación del problema inverso}
\end{table}

\begin{knowledgebox}{Generalidad del problema inverso}
Los problemas inversos no se limitan al dominio de las imágenes. En áreas como el diseño molecular, el plegamiento de proteínas y la planificación robótica, cualquier problema que implique ``inferir información completa desde observaciones parciales'' puede modelarse como un problema inverso. Si puedes definir un mapeo hacia adelante $\mathcal{A}$ razonable para un nuevo dominio, podrás utilizar las técnicas presentadas en esta clase para resolverlo.
\end{knowledgebox}

\subsection{Resumen de la sección}

El problema inverso proporciona un marco unificado que conecta la guía de generación durante la inferencia con los clásicos problemas de procesamiento de señales/imagen computacional. Los elementos clave son: el mapeo hacia adelante conocido $\mathcal{A}$, el valor de observación $y$ y el modelo de generación incondicional preentrenado. A continuación, derivaremos cómo utilizar la fórmula de Bayes para transformar la solución del problema inverso en una modificación de funciones puntuales (score functions).

\section{Fórmula de Bayes y descomposición del score condicional}

La herramienta matemática central para resolver problemas inversos es la \textbf{fórmula de Bayes}, que descompone el score condicional en la suma del score marginal y el score de verosimilitud.

\subsection{Descomposición bayesiana del score condicional}

Nuestro objetivo es calcular el score condicional $\nabla_{x_t} \log p(x_t | y)$ para guiar el proceso de generación durante la eliminación de ruido. Utilizando la fórmula de Bayes:

$$
p(x_t | y) = \frac{p(y | x_t) \cdot p(x_t)}{p(y)}
$$

Tomando logaritmos en ambos lados y calculando el gradiente:

$$
\nabla_{x_t} \log p(x_t | y) = \underbrace{\nabla_{x_t} \log p(x_t)}_{\text{Score marginal (conocido)}} + \underbrace{\nabla_{x_t} \log p(y | x_t)}_{\text{Score de verosimilitud (por calcular)}}
$$

\begin{importantbox}{Score condicional = Score marginal + Score de verosimilitud}
\begin{itemize}
    \item \textbf{Score marginal} $\nabla_{x_t} \log p(x_t)$: esta es exactamente la función de score que los modelos de difusión preentrenados ya aprendieron a predecir y se puede obtener directamente del modelo.
    \item \textbf{Score de verosimilitud} $\nabla_{x_t} \log p(y | x_t)$: mide qué tan probable es que la observación $y$ aparezca dado el estado intermedio actual $x_t$. Esta es la parte que requiere un cálculo adicional (o aproximado).
\end{itemize}
Por lo tanto, todo el problema se reduce a: \textbf{¿cómo aproximar el score de verosimilitud sin entrenar una red neuronal adicional?}
\end{importantbox}

\subsection{Comparación con Classifier Guidance}

\begin{knowledgebox}{Métodos existentes bajo un marco unificado}
Las técnicas de guiado estudiadas anteriormente pueden entenderse de manera unificada bajo este marco:
\begin{itemize}
    \item \textbf{Classifier Guidance}: estimando directamente el score de verosimilitud $\nabla_{x_t} \log p(y | x_t)$ entrenando una red neuronal de clasificación adicional.
    \item \textbf{Classifier-Free Guidance}: aprendiendo conjuntamente con la misma red tanto el ruido condicional como el incondicional, calculando implícitamente el score de verosimilitud a través de la diferencia entre ambos.
\end{itemize}
Ambos métodos requieren entrenamiento. Lo que nosotros vamos a hacer es \textbf{aproximar el score de verosimilitud sin necesidad de entrenamiento}.
\end{knowledgebox}

\subsection{Resumen de la sección}

La descomposición bayesiana nos indica que la generación condicional solo requiere añadir un score de verosimilitud adicional sobre el score marginal del modelo preentrenado. El desafío central restante es: ¿cómo aproximar este score de verosimilitud de manera eficiente?



\section{Aproximación de la función de verosimilitud: de lo ideal a lo factible}

El cálculo directo de $\nabla_{x_t} \log p(y | x_t)$ no es factible, ya que $p(y | x_t)$ involucra un mapeo complejo desde los datos intermedios con ruido $x_t$ hasta la observación $y$. Necesitamos una serie de aproximaciones para hacer viable el cálculo.

\subsection{Dificultades clave: $x_t$ vs. $x_0$}

Podemos escribir directamente la expresión para $p(y | x_0)$ —cuando $\mathcal{A}$ es una función lineal, $y | x_0 \sim \mathcal{N}(\mathcal{A} x_0, \sigma^2 I)$, lo cual constituye una distribución gaussiana sencilla.

Pero lo que necesitamos es $p(y | x_t)$, no $p(y | x_0)$. El problema radica en: \textbf{¿cómo establecer la conexión entre $x_t$ y $x_0$?}

\begin{warningbox}{No se puede sustituir directamente $p(y|x_t)$ por $p(y|x_0)$}
La aproximación más tosca consiste en asumir que $p(y|x_t) \approx p(y|x_0)$, es decir, tratar los datos intermedios con ruido como si fueran datos limpios para calcular. Esta es una aproximación terrible: si $p(x_0|x_t)$ fuera realmente una distribución delta (es decir, si se pudiera recuperar $x_0$ de manera determinista a partir de $x_t$), no tendríamos necesidad del proceso iterativo de eliminación de ruido en absoluto.
\end{warningbox}

\subsection{Fórmula de Tweedie: conexión entre $x_t$ y $x_0$}

La \textbf{fórmula de Tweedie} proporciona un método para estimar la media de $x_0$ a partir de $x_t$:

$$
\hat{x}_0 = \mathbb{E}[x_0 | x_t] = \frac{1}{\sqrt{\bar{\alpha}_t}} \left( x_t + (1 - \bar{\alpha}_t) \nabla_{x_t} \log p(x_t) \right)
$$

Donde $\nabla_{x_t} \log p(x_t)$ es exactamente el score predicho por el modelo preentrenado (o equivalentemente, obtenido a través de la predicción del ruido $\epsilon_\theta(x_t, t)$).

\begin{importantbox}{Función central de la fórmula de Tweedie}
La fórmula de Tweedie nos indica que, dado $x_t$, podemos calcular el \textbf{valor esperado posterior} de $x_0$, denotado como $\hat{x}_0$ (también llamado estimador de eliminación de ruido). Aunque esto representa solo la media y no la distribución completa, sirve como un puente clave para aproximar la función de verosimilitud.
\end{importantbox}

\subsection{Aproximación gaussiana: simplificación de $p(x_0 | x_t)$}

Para hacer viable el cálculo, hacemos una suposición aproximativa (aunque tosca):

$$
p(x_0 | x_t) \approx \mathcal{N}(\hat{x}_0, r_t^2 I)
$$

Es decir, aproximar la distribución posterior de $x_0$ dado $x_t$ como una distribución gaussiana con media dada por la estimación de Tweedie $\hat{x}_0$ y varianza escalar $r_t^2$.

\begin{warningbox}{Limitaciones de la aproximación gaussiana}
Esta aproximación es bastante tosca. Si $p(x_0|x_t)$ fuera realmente una distribución gaussiana, no necesitaríamos el proceso iterativo de eliminación de ruido; bastaría con muestrear $x_0$ directamente desde $x_t$. La verdadera $p(x_0|x_t)$ es una distribución multimodal mucho más compleja que la gaussiana. Sin embargo, sin entrenamiento adicional, esta es una de las mejores aproximaciones posibles.
\end{warningbox}

\subsection{Resumen de la sección}

A través de la fórmula de Tweedie y la aproximación gaussiana, hemos establecido un puente desde $x_t$ hasta $x_0$, transformando la $p(y|x_t)$ que no puede calcularse directamente en una distribución gaussiana susceptible de cálculo analítico. A continuación, derivaremos el algoritmo de guía específico basado en esto.



\section{Solución analítica de problemas inversos lineales: guía por pseudoínverson}

Cuando el mapeo hacia adelante $\mathcal{A}$ es una función lineal (lo cual abarca una gran cantidad de aplicaciones prácticas como desenfocamiento, superresolución, restauración y reconstrucción TC), podemos obtener una expresión analítica para la función de verosimilitud.

\subsection{Configuración del modelo hacia adelante lineal}

Supongamos que $\mathcal{A}$ es lineal, es decir, $y = A x_0 + n$, donde $A$ es una matriz y $n \sim \mathcal{N}(0, \sigma_y^2 I)$ es el ruido de observación. En este caso:

$$
p(y | x_0) = \mathcal{N}(y; A x_0, \sigma_y^2 I)
$$

\begin{knowledgebox}{Matriz $A$ en problemas inversos lineales comunes}
\begin{itemize}
    \item \textbf{Desenfocamiento gaussiano}: $A$ es la matriz correspondiente al núcleo de convolución gaussiana.
    \item \textbf{Superresolución}: $A$ es la matriz de muestreo inferior (downsampling).
    \item \textbf{Restauración de imágenes}: $A$ es una matriz diagonal de máscara (las regiones conocidas son 1, las desconocidas son 0).
    \item \textbf{Reconstrucción TC/MRI}: $A$ es la matriz de proyección/muestreo de Fourier.
\end{itemize}
\end{knowledgebox}

\subsection{Derivación del prior gaussiano conjugado}

Utilizando las propiedades de conjugación de la distribución gaussiana, cuando el prior $p(x_0 | x_t)$ y la verosimilitud $p(y|x_0)$ son ambos distribuciones gaussianas, la distribución marginal $p(y | x_t)$ también es una distribución gaussiana. Específicamente:

\textbf{Prior} (después de aproximación gaussiana):
$$
p(x_0 | x_t) \approx \mathcal{N}(\hat{x}_0, r_t^2 I)
$$

\textbf{Verosimilitud}:
$$
p(y | x_0) = \mathcal{N}(A x_0, \sigma_y^2 I)
$$

\textbf{Marginal} (después de integración sobre $x_0$):
$$
p(y | x_t) = \int p(y | x_0) p(x_0 | x_t) \, dx_0 = \mathcal{N}(A \hat{x}_0, \; r_t^2 A A^\top + \sigma_y^2 I)
$$

\begin{importantbox}{Fórmula de guía por pseudoínverson (Pseudo-Inverse Guidance)}
Basado en la distribución gaussiana marginal anterior, el gradiente de verosimilitud puede calcularse analíticamente:
$$
\nabla_{x_t} \log p(y | x_t) \approx -\frac{1}{2} \nabla_{x_t} \left\| y - A \hat{x}_0 \right\|^2_{(r_t^2 A A^\top + \sigma_y^2 I)^{-1}}
$$
donde $\hat{x}_0$ es la estimación de eliminación de ruido calculada a partir de $x_t$ mediante la fórmula de Tweedie.

Cuando no hay ruido en la observación ($\sigma_y = 0$), la fórmula se simplifica a requerir el cálculo de la \textbf{pseudoínverson} de $A$, por lo que se denomina guía por pseudoínverson.
\end{importantbox}

\begin{warningbox}{La guía por pseudoínverson solo aplica a $\mathcal{A}$ lineal}
La derivación de esta fórmula analítica depende de la linealidad de $\mathcal{A}$. Para mapeos hacia adelante no lineales (como el renderizado que involucra redes neuronales o simulaciones físicas complejas), se deben utilizar métodos más generales (aunque aproximaciones más gruesas).
\end{warningbox}

\subsection{Resumen de la sección}

Para problemas inversos lineales, mediante las propiedades del prior gaussiano conjugado, podemos obtener una expresión analítica para el gradiente de verosimilitud. Este método requiere calcular la pseudoínverson de la matriz de mapeo hacia adelante, por lo que se denomina guía por pseudoínverson. Proporciona una solución teóricamente más rigurosa para tareas clásicas de procesamiento de imágenes como desenfocamiento, superresolución y restauración.
\section{Muestreo de la posterior difusiva: DPS}

Cuando estamos dispuestos a aceptar una aproximación más gruesa a cambio de una aplicabilidad más amplia, obtenemos el método \textbf{DPS (Diffusion Posterior Sampling)}.

\subsection{Aproximación Delta más gruesa}

DPS utiliza una aproximación más extrema que la gaussiana: aproxima directamente $p(x_0 | x_t)$ como una distribución Delta:

$$
p(x_0 | x_t) \approx \delta(x_0 - \hat{x}_0)
$$

Donde $\hat{x}_0$ es la estimación de Tweedie. Esto significa que asumimos que $x_0$ puede recuperarse determinísticamente a partir de $x_t$.

\begin{warningbox}{Costos y beneficios de la aproximación Delta}
\textbf{Costo}: Es una aproximación peor que la gaussiana. Si $p(x_0|x_t)$ realmente fuera una distribución Delta, no sería necesario realizar un proceso iterativo de eliminación del ruido.

\textbf{Beneficio}: El cálculo se vuelve extremadamente simple y \textbf{ya no es necesario que $\mathcal{A}$ sea lineal} —se puede utilizar cualquier mapa hacia adelante diferenciable.
\end{warningbox}

\subsection{Puntuación de verosimilitud del DPS}

Bajo la aproximación Delta:

$$
p(y | x_t) \approx p(y | x_0 = \hat{x}_0) = \mathcal{N}(y; \mathcal{A}(\hat{x}_0), \sigma_y^2 I)
$$

\begin{importantbox}{Fórmula de guía del DPS}
La puntuación de verosimilitud del DPS es:
$$
\nabla_{x_t} \log p(y | x_t) \approx -\frac{1}{\sigma_y^2} \nabla_{x_t} \left\| y - \mathcal{A}(\hat{x}_0) \right\|_2^2
$$
Donde $\hat{x}_0 = \hat{x}_0(x_t)$ es la estimación de eliminación del ruido calculada a partir de $x_t$ mediante la fórmula de Tweedie, y $\nabla_{x_t}$ representa el gradiente respecto a $x_t$ (que requiere retropropagación a través de la dependencia de $\hat{x}_0$ en $x_t$).

\textbf{Significado intuitivo}: Comparar la diferencia entre el resultado del mapa hacia adelante aplicado a la estimación de Tweedie, $\mathcal{A}(\hat{x}_0)$, y la observación real $y$, calcular el gradiente de la pérdida L2 respecto a $x_t$ para servir como señal de guía.
\end{importantbox}

\subsection{Flujo del algoritmo DPS}

La implementación del DPS es extremadamente concisa: solo se requiere agregar una línea de código al ciclo de muestreo estándar de DDPM:

\begin{lstlisting}[caption={Pseudocódigo de guía DPS}]
# Bucle estándar de muestreo DDPM
for t in reversed(range(T)):
    # 1. Paso estándar de eliminación del ruido
    eps_pred = model(x_t, t)           # predecir el ruido
    x0_hat = tweedie_estimate(x_t, eps_pred, t)  # Tweedie
    x_t_minus_1 = ddpm_step(x_t, eps_pred, t)    # paso estándar

    # 2. Guía DPS (¡solo una línea adicional!)
    loss = torch.norm(y - A(x0_hat)) ** 2
    grad = torch.autograd.grad(loss, x_t)[0]
    x_t_minus_1 = x_t_minus_1 - step_size * grad
\end{lstlisting}

\begin{importantbox}{Dos implementaciones equivalentes del DPS}
Bajo el marco de DDPM, las siguientes dos implementaciones son equivalentes:
\begin{enumerate}
    \item Modificar la predicción de ruido $\epsilon$ y luego ejecutar el paso de eliminación del ruido.
    \item Ejecutar primero el paso estándar de eliminación del ruido para obtener $x_{t-1}'$, y luego agregar directamente el término de guía en $x_{t-1}'$.
\end{enumerate}
Esto se debe a que $x_{t-1}$ es una función afín de $\epsilon$, por lo que ambas formas son matemáticamente completamente equivalentes. La segunda implementación es más intuitiva y concisa.
\end{importantbox}

\subsection{Alcance del DPS}

\begin{knowledgebox}{El DPS soporta mapas hacia adelante no lineales}
Una ventaja clave del DPS es que \textbf{no requiere} que $\mathcal{A}$ sea lineal. Siempre que $\mathcal{A}$ sea diferenciable (para permitir la retropropagación y el cálculo de gradientes), se puede aplicar el DPS. Esto permite al DPS manejar una gama más amplia de problemas inversos, tales como:
\begin{itemize}
    \item Recuperación de imágenes con distorsión no lineal
    \item Reconstrucción 3D que involucra renderizadores basados en redes neuronales
    \item Recuperación de señales basada en modelos físicos no lineales
\end{itemize}
El costo es una precisión de aproximación menor (aproximación Delta), pero en la práctica todavía produce resultados muy buenos.
\end{knowledgebox}

\subsection{Resumen de la sección}

DPS convierte el cálculo de la puntuación de verosimilitud en un gradiente de una pérdida L2 simple al adoptar la aproximación Delta más sencilla. Su implementación solo requiere agregar un código de actualización de gradiente adicional al ciclo de muestreo estándar. Aunque la aproximación es más gruesa, esto a cambio soporta mapas hacia adelante no lineales, ampliando significativamente el alcance de aplicación.
\section{Comparación de métodos y compensaciones de diseño}

\subsection{Comparación de tres aproximaciones a la función de verosimilitud}

\begin{table}[H]
\centering
\begin{tabular}{p{0.2\textwidth}p{0.25\textwidth}p{0.2\textwidth}p{0.2\textwidth}}
\toprule
\textbf{Método} & \textbf{Aproximación de $p(x_0|x_t)$} & \textbf{Requisito para $\mathcal{A}$} & \textbf{Precisión de la aproximación} \\
\midrule
Aproximación más rudimentaria & $p(y|x_t) = p(y|x_0)$ & Cualquiera & Peor \\
DPS & $\delta(x_0 - \hat{x}_0)$ & Diferenciable es suficiente & Inferior \\
Pseudo-inversa guiada & $\mathcal{N}(\hat{x}_0, r_t^2 I)$ & Debe ser lineal & Relativamente mejor \\
\bottomrule
\end{tabular}
\caption{Comparación de tres métodos de aproximación a la función de verosimilitud}
\end{table}

\begin{importantbox}{Compensación entre precisión de la aproximación y alcance de aplicación}
Estos tres métodos ilustran una compensación de diseño central:
\begin{itemize}
    \item Las aproximaciones más precisas (como la pseudo-inversa guiada) requieren supuestos más fuertes ($\mathcal{A}$ debe ser lineal), pero son teóricamente más confiables.
    \item Las aproximaciones más rudimentarias (como DPS) tienen requisitos menores para $\mathcal{A}$ (solo necesita ser diferenciable), tienen un alcance de aplicación más amplio, pero la error de aproximación es mayor.
\end{itemize}
No existe ningún método que sea óptimo en todas las circunstancias; la elección depende del caso de uso específico y de las propiedades del mapeo hacia adelante.
\end{importantbox}

\subsection{Limitaciones comunes de todos los métodos}

\begin{warningbox}{Limitaciones inherentes de la guía por aproximación}
Independientemente de la aproximación utilizada, estos métodos de guía durante la inferencia comparten las siguientes limitaciones:
\begin{enumerate}
    \item \textbf{No garantizan convergencia a la distribución posterior correcta}: Todas las aproximaciones introducen errores; los resultados generados pueden no satisfacer completamente las condiciones de restricción.
    \item \textbf{Requieren ajuste de hiperparámetros}: La intensidad de la guía (step size / escala de guía) necesita ser ajustada manualmente, y diferentes tareas y conjuntos de datos pueden requerir valores distintos.
    \item \textbf{Costo computacional}: DPS requiere retropropagación del estimador de Tweedie, lo que incrementa el costo computacional por paso.
    \item \textbf{Las aproximaciones en pasos de tiempo con alto ruido son peores}: En las etapas tempranas (con alto ruido), la calidad de la estimación de $\hat{x}_0$ es muy pobre, lo que hace que la señal de guía sea poco confiable.
\end{enumerate}
\end{warningbox}

\subsection{Resumen de la sección}

Diferentes métodos de aproximación a la función de verosimilitud forman un espectro entre precisión y alcance de aplicación. En aplicaciones prácticas, es necesario seleccionar el método adecuado basándose en las propiedades del mapeo hacia adelante, el presupuesto computacional y los requisitos de calidad de resultados.


\section{Detalles de implementación de la guía bajo el marco DDPM}

Para ayudar a los lectores a comprender cómo funciona DPS a nivel de código, esta sección desglosa detalladamente cómo se combinan los pasos de desruido de DDPM con el término de guía.

\subsection{Repaso del proceso inverso de DDPM}

El proceso inverso (reverse process) de DDPM define cómo muestrear $x_{t-1}$ a partir de $x_t$:

$$
x_{t-1} = \frac{1}{\sqrt{\alpha_t}} \left( x_t - \frac{1-\alpha_t}{\sqrt{1-\bar{\alpha}_t}} \epsilon_\theta(x_t, t) \right) + \sigma_t z
$$

Donde $z \sim \mathcal{N}(0, I)$ es un término de ruido aleatorio, $\sigma_t$ es la desviación estándar del ruido correspondiente al paso de tiempo, y $\epsilon_\theta$ es la red preentrenada para predecir el ruido.

Esta fórmula puede reescribirse en términos de la forma del estimador Tweedie $\hat{x}_0$:

$$
x_{t-1} = c_1(t) \cdot x_t + c_2(t) \cdot \hat{x}_0 + \sigma_t z
$$

Donde $c_1(t)$ y $c_2(t)$ son coeficientes que dependen únicamente del cronograma de ruido (noise schedule). El punto clave es: \textbf{$x_{t-1}$ es una función afín de $\hat{x}_0$, y $\hat{x}_0$ es a su vez una función de $x_t$}.

\subsection{Método de inyección del término de guía}

La operación central de la guía consiste en añadir un término de corrección al $x_{t-1}$ de cada paso:

$$
x_{t-1}^{\text{guided}} = x_{t-1}^{\text{unguided}} - \zeta_t \cdot \nabla_{x_t} \left\| y - \mathcal{A}(\hat{x}_0) \right\|_2^2
$$

Donde $\zeta_t$ es el tamaño de paso (step size) asociado al paso de tiempo.

\begin{warningbox}{El cálculo del gradiente requiere retropropagación}
Nota que el gradiente en $\nabla_{x_t} \| y - \mathcal{A}(\hat{x}_0) \|_2^2$ se calcula con respecto a $x_t$, mientras que $\hat{x}_0$ depende de $x_t$ (a través de la fórmula Tweedie y la red de predicción de ruido $\epsilon_\theta$). Por lo tanto, calcular este gradiente implica realizar una retropropagación a través de $\epsilon_\theta$, lo que significa:
\begin{enumerate}
    \item Cada paso de guía requiere una retropropagación adicional, con un costo computacional aproximadamente 2-3 veces mayor que la propagación hacia adelante.
    \item Los parámetros de la red de predicción de ruido no se actualizan; simplemente utilizamos la información del gradiente para modificar la trayectoria de muestreo.
    \item Es necesario habilitar el rastreo de gradientes en PyTorch mediante \texttt{x\_t.requires\_grad\_(True)}.
\end{enumerate}
\end{warningbox}

\subsection{Consideraciones prácticas sobre el cronograma del tamaño de paso}

\begin{knowledgebox}{Estrategia de selección del tamaño de paso $\zeta_t$}
La elección del tamaño de guía $\zeta_t$ es crucial para la calidad del resultado:
\begin{itemize}
    \item \textbf{Demasiado grande}: Los resultados generados cumplen estrictamente con la restricción $\mathcal{A}(x_0) = y$, pero pueden desviarse de la variedad (manifold) de los datos, generando artefactos.
    \item \textbf{Demasiado pequeño}: Los resultados son realistas pero no satisfacen suficientemente las condiciones de restricción.
    \item \textbf{Estrategias comunes}: Utilizar un tamaño de paso constante o una atenuación lineal según el paso de tiempo (tamaño de paso pequeño al principio, ya que la estimación de $\hat{x}_0$ es imprecisa; tamaño de paso grande al final, ya que la estimación es más precisa).
    \item \textbf{Estrategias adaptativas}: Ajustar dinámicamente el tamaño de paso según el valor de $\| y - \mathcal{A}(\hat{x}_0) \|$.
\end{itemize}
En la práctica, generalmente se requiere realizar una búsqueda en cuadrícula sobre tareas específicas para determinar el tamaño de paso óptimo.
\end{knowledgebox}

\subsection{Resumen de la sección}

Bajo el marco DDPM, la implementación de la guía DPS se reduce a añadir un término de corrección de gradiente después de cada paso de desruido. Los detalles de implementación clave incluyen la retropropagación del gradiente, la selección del tamaño de paso y la gestión del costo computacional.

\section{Intuición comprensiva guiada durante la inferencia}

\subsection{Visión geométrica de las trayectorias curvas}

Entendido desde la intuición geométrica, la esencia de la guía durante la inferencia es el \textbf{ruido que se desvía en una trayectoria curva}:

\begin{enumerate}
    \item Sin guía, partiendo de $x_T$ (ruido puro), el modelo difusivo sigue la dirección indicada por las puntuaciones marginales $\nabla_{x_t} \log p(x_t)$, dirigiendo la trayectoria hacia algún punto $x_0$ en el subespacio de datos.
    \item Al agregar guía, la puntuación de verosimilitud $\nabla_{x_t} \log p(y|x_t)$ aplica una fuerza adicional a la trayectoria en cada paso, haciéndola inclinarse hacia aquellos puntos de datos que satisfacen la condición $\mathcal{A}(x_0) \approx y$.
    \item Finalmente, el $x_0$ al que llega la trayectoria es tanto un punto razonable en el subespacio de datos (garantizado por las puntuaciones marginales) como uno que cumple con las restricciones condicionales (guiado por las puntuaciones de verosimilitud).
\end{enumerate}

\begin{knowledgebox}{¿Por qué los modelos difusivos son especialmente adecuados para la guía durante la inferencia?}
Los GAN y VAE son modelos de generación en un solo paso: el mapeo de ruido a datos es una única propagación hacia adelante. Una vez que este mapeo se fija, resulta difícil aplicar controles adicionales durante la inferencia.

En cambio, el proceso generativo de los modelos difusivos/flujo consiste en una serie de \textbf{pequeñas iteraciones}, donde cada paso puede ser ajustado finamente. Esto es similar a la navegación: si llegas al destino en un solo paso, solo puedes ir en línea recta; pero si lo divides en muchos pasos, puedes ajustar la dirección en cada uno según nueva información —esa es precisamente la esencia de la guía durante la inferencia.
\end{knowledgebox}

\subsection{Resumen de la sección}

La esencia geométrica de la guía durante la inferencia consiste en aplicar una fuerza correctora direccional adicional en cada paso de la trayectoria de desruido, haciendo que el resultado generado final satisfaga simultáneamente los objetivos de ser ``realista'' y ``cumplir con las condiciones''. Las características iterativas del modelo difusivo lo hacen naturalmente adecuado para este tipo de operación.
\section{Síntesis y ampliación}

\subsection{Guía de escenarios aplicables para cada método}

\begin{table}[H]
\centering
\begin{tabular}{p{0.22\textwidth}p{0.35\textwidth}p{0.33\textwidth}}
\toprule
\textbf{Escenario} & \textbf{Método recomendado} & \textbf{Justificación} \\
\midrule
Reparación de imágenes & Manipulación del Score (RePaint) o DPS & La operación con máscara es sencilla; ambas clases de métodos son viables \\
Superresolución & Guiado por pseudoinversa & El submuestreo es una operación lineal, por lo que se puede utilizar la solución analítica \\
Desenfoque gaussiano & Guiado por pseudoinversa & La convolución es una operación lineal \\
Reconstrucción de CT/MRI & Guiado por pseudoinversa & La operación de proyección es lineal \\
Recuperación de distorsión no lineal & DPS & Se requiere manejar el operador no lineal $\mathcal{A}$ \\
Generación con control de layout & Manipulación del Attention Map & Es necesario controlar las relaciones espaciales correspondientes \\
Transferencia de estilo & Manipulación del Score & Implica el intercambio de características entre imágenes \\
\bottomrule
\end{tabular}
\caption{Guía para la selección de métodos de guiado durante la inferencia en diferentes escenarios de aplicación}
\end{table}

\subsection{Repaso del contenido central de esta clase}

Esta clase presenta sistemáticamente el marco y los métodos básicos de \textbf{generación guiada durante la inferencia}, abarcando los siguientes contenidos centrales:

\begin{enumerate}
    \item \textbf{Motivación}: Aprovechar las características de generación iterativa de los modelos de difusión preentrenados para inyectar señales de guía durante la inferencia, logrando así una generación condicional sin necesidad de volver a entrenar.
    \item \textbf{Manipulación del Score}: Métodos prácticos que manipulan directamente el score de des ruido o los mapas de atención (como RePaint), sin garantías teóricas pero ampliamente efectivos.
    \item \textbf{Marco de problemas inversos}: Modelar la generación condicional como un problema inverso $y = \mathcal{A}(x) + n$, proporcionando herramientas de análisis sistemáticas.
    \item \textbf{Descomposición bayesiana}: $\nabla_{x_t} \log p(x_t|y) = \nabla_{x_t} \log p(x_t) + \nabla_{x_t} \log p(y|x_t)$, reduciendo el problema a la aproximación del score de verosimilitud.
    \item \textbf{Guiado por pseudoinversa}: Basado en una aproximación gaussiana, aplicable a mapeos hacia adelante lineales con una precisión de aproximación relativamente buena.
    \item \textbf{DPS}: Basado en una aproximación Delta, aplicable a cualquier mapeo hacia adelante diferenciable; su implementación es extremadamente simple pero la aproximación es más rugosa.
\end{enumerate}

\subsection{Conexión con la siguiente clase}

Esta clase cubre las dos primeras de tres categorías principales de tecnologías de guiado durante la inferencia. La próxima clase (Lecture 12: Inference-Time Guided Generation 2) discutirá tecnologías más avanzadas de la tercera categoría, que podrían incluir:
\begin{itemize}
    \item Métodos de aproximación del score de verosimilitud más precisos
    \item Guiado durante la inferencia basado en optimización (como DDRM, MCG, etc.)
    \item Aplicaciones específicas del guiado durante la inferencia en dominios como la generación 3D y de video
\end{itemize}

\subsection{Lecturas adicionales}

\begin{itemize}
    \item \textbf{RePaint}: Lugmayr et al., ``RePaint: Inpainting using Denoising Diffusion Probabilistic Models,'' CVPR 2022.
    \item \textbf{DPS}: Chung et al., ``Diffusion Posterior Sampling for General Noisy Inverse Problems,'' ICLR 2023.
    \item \textbf{DDRM}: Kawar et al., ``Denoising Diffusion Restoration Models,'' NeurIPS 2022.
    \item \textbf{Classifier Guidance}: Dhariwal \& Nichol, ``Diffusion Models Beat GANs on Image Synthesis,'' NeurIPS 2021.
    \item \textbf{Sobre problemas inversos y modelos de difusión}: El orador recomienda consultar el paper de revisión (survey) para conocer más sobre técnicas de manipulación del score y resolución de problemas inversos.
\end{itemize}

%% --- Fin del contenido --- %%

\end{document}
